% Livro C — um exercício que começa numa página e termina na seguinte (prompt 03, §52).
% O resultado obrigatório é UM exercício com DUAS âncoras, não dois exercícios.
\documentclass[11pt,oneside]{book}
% Preâmbulo comum das fixtures do scan (D48). Fontes Latin Modern com T1, para que a camada de
% texto traga os acentos como caracteres, e babel em português, para que \chapter imprima
% "Capítulo" como um livro brasileiro imprime.
\usepackage[T1]{fontenc}
\usepackage[utf8]{inputenc}
\usepackage[brazil]{babel}
\usepackage{lmodern}
\usepackage{amsmath,amssymb}
\usepackage{tikz}
\usepackage[a4paper,margin=2.5cm]{geometry}
\usepackage{enumitem}
\setlist[enumerate,2]{label=\alph*)}
\makeatletter\@ifundefined{c@chapter}{\newcounter{exemplo}}{\newcounter{exemplo}[chapter]}\makeatother
\newenvironment{exemplo}{\par\medskip\refstepcounter{exemplo}\noindent\textbf{Exemplo \theexemplo.}\ }{\par\medskip}
\newcommand{\blocoexercicios}[1][Exercícios]{\section*{#1}}

\begin{document}
\chapter{Sequências}
\section{Progressões}
Uma progressão aritmética é uma sequência em que a diferença entre termos consecutivos é
constante.

\blocoexercicios
\begin{enumerate}[label=\textbf{\arabic*.}]
\item Calcule o décimo termo da progressão $2, 5, 8, \ldots$
\vspace*{19.5cm}
\item Considere a progressão aritmética de primeiro termo $a_1 = 7$ e razão $r = 4$.
Este enunciado é longo de propósito, para atravessar o pé da página e continuar na seguinte.
Determine o termo geral da progressão, justificando cada passagem com cuidado, e depois
responda às perguntas abaixo, sempre mostrando as contas que levam a cada resposta.
  \begin{enumerate}
  \item Qual é o vigésimo termo?
  \item Qual é a soma dos dez primeiros termos?
  \item Algum termo da progressão é igual a $100$? Justifique.
  \item Qual é o menor termo maior que $500$?
  \item Quantos termos há entre $50$ e $300$?
  \item Escreva a progressão em ordem decrescente a partir do décimo termo.
  \end{enumerate}
\item Calcule a razão da progressão $10, 7, 4, \ldots$
\end{enumerate}
\end{document}
